\documentclass{exam}
\usepackage{../../shah311}
\usepackage[linguistics]{forest}

\hypersetup{colorlinks=true, linktoc=section, linkcolor=blue}

\pagestyle{headandfoot}
\firstpageheadrule
\runningheadrule
\firstpageheader{Name: Jeevan Shah}{Math 477 Homework 3}{RU ID: 242001315 \\ netID: jds508}
\runningheader{Name: Jeevan Shah}{Math 477 Homework 3}{RU ID: 242001315 \\ netID: jds508}
\firstpagefooter{}{}{}
\runningfooter{ }{\thepage}{ }

\printanswers

\begin{document}
\colorbox{red}{\underline{1:}} Ten percent of a company's life insurance policyholders are smokers. The rest are non-smokers. For each non-smoker, the probability of dying during the year is 0.01. For each smoker, the probability of dying during the year is 0.05. Given that a policyholder died, what is the probability that the policyholder was a smoker?
\begin{solution}
    Let $S$ be the event that a policyholder is a smoker and $D$ the event that a policyholder dies during the year.
    \begin{center}
        \begin{forest}
            [
                [
                    $S \left(0.1\right)$
                    [
                        $D \left(0.05\right)$
                    ]
                    [
                        $D^c \left(0.95\right)$
                    ]
                ]
                [
                    $S^c \left(0.9\right)$
                    [
                        $D \left(0.01\right)$
                    ]
                    [
                        $D^c \left(0.99\right)$
                    ]
                ]
            ]
        \end{forest}
    \end{center}
    By Bayes' formula we have
    \begin{align*}
        P(S \mid D) &= \frac{P(D \mid S)P(S)}{P(D \mid S)P(S) + P(D \mid S^c)P(S^c)} \\
        &= \frac{(0.05)(0.1)}{(0.05)(0.1) + (0.01)(0.9)} \\
        &= \frac{0.005}{0.014} \\
        &= \boxed{\frac{5}{14} \approx 0.357}
    \end{align*}
\end{solution}

\colorbox{red}{\underline{2:}} An actuary studying the insurance preferences of automobile owners makes the following conclusions:
\begin{enumerate}
    \item[(i)] An automobile owner is twice as likely to purchase collision coverage as disability coverage.
    \item[(ii)] The event that an automobile owner purchases collision coverage is independent of the event that he or she purchases disability coverage.
    \item[(iii)] The probability that an automobile owner purchases both collision and disability coverages is 0.15.
\end{enumerate}
Calculate the probability that an automobile owner purchases neither collision nor disability coverage.

\begin{oneparchoices}
    \choice 0.18
    \CorrectChoice 0.33
    \choice 0.48
    \choice 0.67
    \choice 0.82
\end{oneparchoices}
\begin{solution}
    Let $C$ be the event that an owner purchases collision coverage and $D$ the event that an owner purchases disability coverage. From (i) we have $P(C) = 2P(D)$, and from (ii) and (iii) we have
    \begin{align*}
        &P(C \cap D) = P(C)P(D) = 0.15 \\
        \Rightarrow& 2P(D)^2 = 0.15 \\
        \Rightarrow& P(D) = \sqrt{0.075} \approx 0.2739 \\
        \Rightarrow& P(C) = 2\sqrt{0.075} \approx 0.5477
    \end{align*}
    So by inclusion-exclusion
    \begin{align*}
        P(C^c \cap D^c) &= 1 - P(C \cup D) \\
        &= 1 - \left(P(C) + P(D) - P(C \cap D)\right) \\
        &\approx 1 - (0.5477 + 0.2739 - 0.15) \\
        &= \boxed{0.3284 \approx 0.33}
    \end{align*}
    and thus \fbox{B} is the correct answer.
\end{solution}
\footnotetext{\LaTeX\, code for this document can be found on github \href{https://github.com/jeevanshah07/MATH477/blob/main/homework/homework3/main.tex}{\underline{here}}}

\end{document}
